\documentclass[11pt]{article}
\usepackage{amsmath,amssymb,booktabs,geometry,siunitx}
\usepackage[hidelinks]{hyperref}
\geometry{margin=1in}
\newcommand{\EqID}[1]{\label{eq:#1}}
\begin{document}

\title{Mathematical formulation}
\author{cse-heat-equation-example}
\date{}
\maketitle

\section{Scope and basis}

This document defines the dimensionless reference case implemented in
\path{src/heat_verification/model.py}. Every quantity is dimensionless.
The project addresses code and solution verification only. It introduces no
material property, experiment, or dimensional conversion.

\section{Continuous problem}

\subsection{EQ-HEAT-001: governing equation}

The project-declared equation is
\begin{equation}
  \frac{\partial \theta}{\partial Fo}
  = \frac{\partial^2 \theta}{\partial \xi^2},
  \qquad 0<\xi<1,\quad Fo>0.
  \EqID{heat-pde}
\end{equation}
Here \(\theta\) is dimensionless temperature, \(\xi\) is dimensionless
position, and \(Fo\) is Fourier number. The equation is one-dimensional,
linear, homogeneous, and applied only to the smooth single-mode case below.

\subsection{BC-HEAT-001: boundary conditions}

The boundaries remain exactly zero:
\begin{equation}
  \theta(0,Fo)=\theta(1,Fo)=0,
  \qquad Fo\geq 0.
  \EqID{heat-bc}
\end{equation}
Both solvers reset the two endpoint values after every time step. The boundary
invariant is protected by \path{tests/test_scientific_behavior.py}.

\subsection{IC-HEAT-001: initial condition}

The initial field is the first boundary-compatible sine mode:
\begin{equation}
  \theta(\xi,0)=\sin(\pi\xi),
  \qquad 0\leq\xi\leq1.
  \EqID{heat-ic}
\end{equation}
Its amplitude is one on the dimensionless temperature basis.

\section{Reference objects}

\subsection{REF-FOURIER-001: continuous reference}

Separation of variables gives the internally derived continuum solution
\begin{equation}
  \theta_{\mathrm{ref}}(\xi,Fo)
  =\exp(-\pi^2 Fo)\sin(\pi\xi).
  \EqID{heat-exact}
\end{equation}
Direct differentiation recovers Eq.~\eqref{eq:heat-pde}; evaluating at the
boundaries and at \(Fo=0\) recovers Eqs.~\eqref{eq:heat-bc} and
\eqref{eq:heat-ic}. The public evaluator is
\path{heat_verification.reference_solution}.

\subsection{REF-SEMIDISCRETE-001: centered-space reference}

For \(N\) uniform intervals, \(\Delta\xi=1/N\), and the first grid sine mode,
the centered second-difference eigenvalue is
\begin{equation}
  \lambda_h
  =-\frac{4}{(\Delta\xi)^2}
    \sin^2\!\left(\frac{\pi}{2N}\right).
  \EqID{semidiscrete-eigenvalue}
\end{equation}
The exact-in-time centered-space field is
\begin{equation}
  \theta_h(\xi_i,Fo)
  =\exp(\lambda_h Fo)\sin(\pi\xi_i).
  \EqID{semidiscrete-reference}
\end{equation}
This project derivation isolates temporal error without conflating it with the
continuous reference. The public evaluator is
\path{heat_verification.semidiscrete_reference}.

\section{Validity and claim boundary}

The references are valid for the declared uniform grid, first sine mode, zero
Dirichlet boundaries, and nonnegative \(Fo\). They do not validate a physical
heat-transfer model. The retained study evaluates \(Fo=0.1\) and records all
grid and time-step conditions with each result.

\end{document}
